\documentclass[10pt,a4paper]{amsart}
\usepackage[margin=19mm]{geometry}
\usepackage{amsmath,amssymb,mathtools}
\usepackage[scr=rsfs,cal=euler]{mathalfa}
\usepackage{xfrac}
\usepackage[unicode,hidelinks,bookmarks=false]{hyperref}
\newcommand{\ie}{i.e.\ }
\newcommand{\cf}{cf.\ }
\newcommand{\resp}{respectively\ }
\newcommand{\Subsection}{\subsection}
\providecommand{\nobreakdash}{\nobreak\hbox{-}}
\setlength{\emergencystretch}{2em}
\multlinegap0pt
\usepackage{bm}
\usepackage{bbm}
\usepackage{dsfont}
\usepackage{xspace}
\usepackage{cancel}
\usepackage{mathtools}
\usepackage{amsthm}
\makeatletter
\@ifundefined{lemma}{\newtheorem{lemma}{Lemma}}{}
\@ifundefined{theorem}{\newtheorem{theorem}{Theorem}}{}
\makeatother
\makeatletter
\expandafter\ifx\csname example\endcsname\relax
\newenvironment{example}[1][]%
{\begin{innerexample}[#1]\pushQED{\qed}}%
{\popQED\end{innerexample}}
\fi
\expandafter\ifx\csname proof\endcsname\relax
\renewenvironment{proof}[1][Proof]{\par\pushQED{\qed}\normalfont\textit{#1.}~}{\popQED\par}
\fi
\makeatother
\title[Half-Global Deadbeat Parking for Dubins Vehicle]{Half-Global Deadbeat Parking for Dubins Vehicle}
\author{Miroslav Krsti\'c, Kwang Hak Kim,, Velimir Todorovski}
\date{Source: arXiv:2509.25571v1 (CC BY 4.0)}
\begin{document}
\maketitle

\noindent\emph{Source and licence.} Half-Global Deadbeat Parking for Dubins Vehicle. Version arXiv:2509.25571v1 (CC BY 4.0), distributed under the Creative Commons Attribution 4.0 International licence, as recorded in the arXiv licence metadata for this exact version and shown on the abstract page https://arxiv.org/abs/2509.25571v1. Full rights notice: licenses/arxiv-2509.25571-CC-BY-4.0.txt.\par\smallskip

\noindent\textit{Editorial scope.} Counted unit: the Theorem whose source label is \texttt{thm:Dubins-FT-stabilize} in the article (source0.tex lines 207--258, source label \texttt{thm:Dubins-FT-stabilize}), delivered together with the article's own complete original proof of it (source0.tex lines 260--365, one \texttt{proof} environment of 106 lines). No mathematics is written, repaired or re-derived: statement and proof are the article's own text line for line. Required context retained uncounted: eq:unicycle\_polar (lines 144--157); eq:unicycle\_polar-delta (lines 146--156); eq:unicycle\_polar-gamma (lines 146--156); eq:unicycle\_polar-rho (lines 146--156); lem1 (lines 177--184); lem2 (lines 190--197). Every definition, display and label of those lines is retained unchanged. Omitted from this package and not redistributed: 1--143, 157--176, 185--189, 198--206, 259--259, 366--748. Nothing inside a retained excerpt is omitted or altered: every retained body is delivered exactly as the article's lines are written, so no omission receipt of any kind accompanies this package. Citations: the retained excerpts carry no citation command at all, so no bibliography entry of the article is transcribed and no reference is redistributed. Reference adaptation: none was needed and none was made; every reference inside the delivered unit resolves inside it, so no unresolved reference or citation is produced. Printed slips kept literal: no printed word, symbol, hypothesis or number of the counted statement or of its proof was altered. Layout and class: the article's own class is \texttt{ieeeconf} with packages including mathrsfs, xcolor, amsmath, amssymb, amsfonts, graphicx, algorithm, algorithmic, hyperref, textcomp, setspace, lipsum, mathtools, enumerate; the wrapper is the lane's pinned \texttt{amsart} editorial wrapper at 10pt on A4, which restates the theorem-like environments the retained text needs and adds the article's own macro definitions that the retained text needs (none), with extra wrapper packages (bm, bbm, dsfont, xspace, cancel, mathtools). The delivered numbering is the unit's own. Assets: no retained span contains a figure environment, a graphics-inclusion command, a PDF-page inclusion, a table or a TikZ picture; no image file is copied into this package, the package declares no asset dependency and no image file is redistributed with it. Licence: version arXiv:2509.25571v1 of this article is distributed under the Creative Commons Attribution 4.0 International licence, as recorded in the arXiv licence metadata for this exact version and shown on the abstract page https://arxiv.org/abs/2509.25571v1. A full rights notice is delivered at licenses/arxiv-2509.25571-CC-BY-4.0.txt. Characters: every retained excerpt is pure ASCII, so the delivered engine, XeLaTeX with its default Latin Modern text font, reports no missing character.
\begin{subequations}
\label{eq:unicycle_polar}
\begin{eqnarray}
\label{eq:unicycle_polar-rho}
\dot{\rho} &=& 
%- k_1 \rho \cos^2(\gamma)
-v \cos\gamma\\
\label{eq:unicycle_polar-delta}
\dot{\delta} &=&  \frac{v}{\rho} \sin\gamma
\\
\label{eq:unicycle_polar-gamma}
\dot{\gamma} &=&  \frac{v}{\rho} \sin\gamma -\omega  \,,
\end{eqnarray} 
\end{subequations}

\begin{lemma}
\label{lem1}
Let $a>0$ and let $V:[0,\rho_0]\to\mathbb{R}_{\ge 0}$ be a continuously differentiable function. Then, for $\rho\in(0,\rho_0]$,
\begin{itemize}
\item $\dfrac{dV}{d\rho} \geq \dfrac{a}{\rho}V$ implies $V(\rho)\leq V(\rho_0) \left(\dfrac{\rho}{\rho_0}\right)^a$. 
\item $\dfrac{dV}{d\rho} \geq \dfrac{a}{\rho^2}V$ implies $V(\rho)\leq V(\rho_0) {\rm e}^{a\left(1/{\rho_0}-{1}/{\rho}\right)}$. 
\end{itemize}
\end{lemma}

\begin{lemma}
\label{lem2}
Consider $\dot\rho = - v\cos\gamma$ for $v>0$, let $\alpha\in\mathcal{K}$, and denote
\begin{equation}
t_1 = \frac{\rho_0}{v}\sqrt{1+\alpha(1)}\,.
\end{equation}
For all $t\in [0,T)$  for which the solution exists, if $\cos\gamma(t)>0$ and $\tan^2\gamma(t)\leq \alpha(\rho(t)/\rho_0)$ for all $t\in[0,T)$, then $\rho(t)\leq \rho_0(1-t/t_1)$ for all $t\in[0,T)$.
\end{lemma}

\begin{theorem}%[Genova CLF] 
\label{thm:Dubins-FT-stabilize}
For $v=\mbox{\rm const}>0$,  consider \eqref{eq:unicycle_polar}. 
Let
\begin{eqnarray}\label{eq-control-omega}
\omega &=& \dfrac{v}{\rho}\bigl\{
\sin\gamma +\cos^2\gamma\bigl[\cos\gamma
(1+c_1c_2)\delta\nonumber \\
&&+(c_1+c_2)\sin\gamma
\bigr]\bigr\}
\end{eqnarray}
and $c_1,c_2 \geq \underline c :=\min\{c_1,c_2\} >1$. For all $\rho_0>0$, $\delta_0\in\mathbb{R}$ and $\gamma_0\in(-\pi/2,\pi/2)$ the following holds:
\begin{equation}
\label{eq-rho-bound}
\rho(t)\leq \rho_0(1-t/t_1)
%- \frac{v}{\sqrt{1+M^2(\delta^2_0 +\tan^2\gamma_0)}}t
\end{equation}
\begin{equation}
\label{eq-deltan-bound}
B^2(t) \leq M^2(c_1)\left( 1- 
t/t_1
%\frac{v}{\rho_0\sqrt{1+M^2(\delta^2_0 +\tan^2\gamma_0)}}t
\right)^{2\underline c}B^2_0
\end{equation}
\begin{align}
\label{eq-omega(t)-bound}
|\omega(t)| \leq & \frac{v}{\rho_0} \,
% \nonumber\\
% &\times 
\sqrt{2}(1+\max\{c_1c_2,c_1+c_2\})\nonumber\\
&\times M(c_1)\left( 1- t/t_1
%\frac{v}{\rho_0\sqrt{1+M^2(\delta^2_0 +\tan^2\gamma_0)}}t
\right)^{\underline c -1}B_0
\end{align}
where $B_0 = B(\delta_0,\gamma_0)$, $B(\delta,\gamma) \coloneqq \sqrt{\delta^2 + \tan^2\gamma}$ and
\begin{equation}
M(s) := 
%\left(
1+\frac{s^2}{2}  + s\sqrt{1+\frac{s^2}{4} }\,, \quad s\geq 0
%\right)^2, 
\label{eq:M}
\end{equation}
for all $t\in\left[0, \min\left\{t_1,T\right\} \right)$, where 
\begin{equation}
t_1(\rho_0,\delta_0, \gamma_0, v,c_1)= \frac{\rho_0}{v}{\sqrt{1+M^2(c_1)B^2_0}}
\end{equation}
and $T$ is the interval of existence of the system's solutions, at which $\rho(T)=0$. In particular, for every $\Omega>0$, the angular velocity $\omega(t)$ is maintained in the interval $[-\Omega,\Omega]$ by choosing the forward velocity such that 
\begin{align}\label{eq-slowdown}
v\leq& \;\Omega\dfrac{\rho_0}{\sqrt{\delta_0^2 +\tan^2\gamma_0}}\nonumber\\
&\times\frac{1}{ \sqrt{2}M(c_1)(1+\max\{c_1c_2,c_1+c_2\})}\,.  
\end{align}
\end{theorem}

\begin{proof}
% Let us start by rewriting \eqref{eq:unicycle_polar-delta} and \eqref{eq:unicycle_polar-gamma} as
% \begin{subequations}
% \label{eq:unicycle_polar-tan}
% \begin{align}
% \label{eq:unicycle_polar-delta-tan}
% \frac{d}{dt}{\delta} &=  \frac{v\cos\gamma}{\rho} \tan\gamma
% \\
% \label{eq:unicycle_polar-gamma-tan}
% \frac{d}{dt}\tan{\gamma} &= \frac{1}{\cos^2\gamma}\left(\frac{v\cos\gamma}{\rho}\tan\gamma -\omega\right)\nonumber\\
% &=: -\frac{v\cos\gamma}{\rho}\bar\omega \,,
% \end{align} 
% \end{subequations}
% where $\omega = \frac{v\cos\gamma}{\rho} \left(\tan\gamma +\cos^2\gamma \, \bar\omega\right)$. 
First, with the backstepping transformation $\zeta = \tan\gamma + c_1 \delta$ and defining $\omega = \frac{v\cos\gamma}{\rho} \left(\tan\gamma +\cos^2\gamma \, \bar\omega\right)$, we rewrite \eqref{eq:unicycle_polar-delta} and \eqref{eq:unicycle_polar-gamma} as
\begin{subequations}
\label{eq:unicycle_polar-zeta}
\begin{align}
\label{eq:unicycle_polar-delta-zeta}
\dot{\delta} &=  \frac{v\cos\gamma}{\rho} \left(-c_1\delta+\zeta\right) %\quad 
% \frac{v\cos\gamma}{\rho}(\zeta - c_1 \delta)
\\
\label{eq:unicycle_polar-gamma-zeta}
\dot\zeta &= 
\frac{v\cos\gamma}{\rho}\left(c_1\tan\gamma -\bar\omega\right) \,.
%\quad 
%\frac{v\cos\gamma}{\rho}(c_1 \tan \gamma + \bar{\omega}) \,.
\end{align} 
\end{subequations}
Introduce the Lyapunov function $V = \delta^2 + \zeta^2$. Its derivative along the solutions of \eqref{eq:unicycle_polar-zeta} is $\dot V = 2 \frac{v\cos\gamma}{\rho}\left[
-c_1\delta^2 +\zeta\left(\delta+c_1\tan\gamma - \bar\omega\right)
\right]$.
% \begin{equation}
% \dot V = 
% % 2 \frac{v\cos\gamma}{\rho}\left[
% % -c_1\delta^2 +\zeta\left(\delta+c_1\tan\gamma - \bar\omega\right)
% % \right] \quad = 
% 2 \frac{v\cos\gamma}{\rho}\left[
% -c_1\delta^2 +\zeta\left(\delta+c_1\tan\gamma - \bar\omega\right)
% \right] \,.
% \end{equation}
Take the control as $\bar\omega = 
%c_2 \zeta+\delta+c_1\tan \gamma \quad =
c_2 \zeta + \delta + c_1\tan \gamma $, which gives \eqref{eq-control-omega} and yields
\begin{equation}
\label{eq-Vdot-closed-loop}
\dot V = -2 \frac{v\cos\gamma}{\rho}\left(
c_1\delta^2 +c_2\zeta^2\right)\,.
\end{equation}
Dividing \eqref{eq-Vdot-closed-loop} by \eqref{eq:unicycle_polar-rho} one obtains
\begin{equation}
\frac{dV}{d\rho} = \frac{2}{\rho}\left(
c_1\delta^2 +c_2\zeta^2\right) \geq 2\underline c \frac{V}{\rho}\,. \label{eq:dVdrho}
\end{equation}
% In order to use the comparison principle for \eqref{eq:dVdrho}, we treat $\tilde{\rho}(t) = \rho_0 - \rho(t)$ as an (increasing) time-like variable.  
% Thus, \eqref{eq:dVdrho} becomes 
% \begin{equation}
% \frac{{\rm d} V}{{\rm d} \tilde{\rho}} \le 2 \underline c \frac{V}{\tilde{\rho} - \rho_0}
% \end{equation}
% Taking into account that $\tilde \rho (0) = 0$, using the comparison principle and substituting $\rho = \rho_0 - \tilde{\rho}$, we get 
From Lemma \ref{lem1}, $V(\rho) \leq \left(\dfrac{\rho}{\rho_0}\right)^{2\underline c} V(\rho_0)$, namely,
\begin{equation}
\label{eq-tandelrho0}
\delta^2+\zeta^2 \leq \left(\frac{\rho}{\rho_0}\right)^{2\underline c} \left(\delta^2_0+\zeta^2_0\right)\,.
\end{equation}
Observing that $\delta^2 + \zeta^2 = \begin{bmatrix}\delta &\tan \gamma\end{bmatrix} P(c_1)  \begin{bmatrix}\delta &\tan \gamma\end{bmatrix}^\top$ where
\begin{equation}
P(c_1) = \begin{bmatrix}
1+c_1^2 &c_1\\c_1 &1
\end{bmatrix},
\end{equation}
and with \eqref{eq-tandelrho0}, one gets  
\begin{equation}
\label{eq-tandelrho}
% \delta^2+\tan^2\gamma \leq \left(\frac{\rho}{\rho_0}\right)^{2\underline c} M^2\left(\delta^2_0+\tan^2\gamma_0\right)\,.
\delta^2+\tan^2\gamma \leq \left(\frac{\rho}{\rho_0}\right)^{2\underline c} M^2(c_1)\left(\delta^2_0+\tan^2\gamma_0\right)\,,
\end{equation}
where $M^2(c_1) = \lambda_{\rm max}(P(c_1))/\lambda_{\rm min}(P(c_1))$ and is defined in \eqref{eq:M}. Inequality \eqref{eq-tandelrho} ensures that 
\begin{equation}
% \delta^2(t)+\tan^2\gamma(t) \leq M^2( \delta^2_0+\tan^2\gamma_0)\,,
%\delta^2(t)+
\tan^2\gamma(t) \leq M^2(c_1)\left(\frac{\rho}{\rho_0}\right)^{2\underline c} ( \delta^2_0+\tan^2\gamma_0)\,,
\end{equation}
which in turn guarantees that $\gamma(t)$ remains in $(-\pi/2,\pi/2)$ and, therefore, $\cos\gamma(t)$ remains positive. 
% and thus $\cos\gamma = 1/\sqrt{1+\tan^2\gamma}$. Consequently, 
% \begin{equation}
% % -\cos\gamma(t) \leq - \frac{1}{\sqrt{1+M^2(\delta^2_0 +\tan^2\gamma_0)}}\,. 
% -\cos\gamma(t) \leq - \frac{1}{\sqrt{1+M^2(c_1)(\delta^2_0 +\tan^2\gamma_0)}}\,. 
% \end{equation}
% Hence
% \begin{equation}
% % \dot\rho = - v\cos\gamma \leq -  \frac{v}{\sqrt{1+M^2(\delta^2_0 +\tan^2\gamma_0)}}\,, 
% \dot\rho = - v\cos\gamma \leq -  \frac{v}{\sqrt{1+M^2(c_1)(\delta^2_0 +\tan^2\gamma_0)}}\,, 
% \end{equation}
% which gives \eqref{eq-rho-bound}. Substituting \eqref{eq-rho-bound} into \eqref{eq-tandelrho}
Using Lemma \ref{lem2}, one gets \eqref{eq-deltan-bound}. From \eqref{eq-control-omega}, 
\begin{equation}
|\omega\rho| \leq v\sqrt{2}(1+\max\{c_1c_2,c_1+c_2\})\sqrt{\delta^2+\tan^2\gamma}\,.
\end{equation}
Substituting \eqref{eq-tandelrho} and with \eqref{eq-rho-bound}, one gets 
\eqref{eq-omega(t)-bound}. Finally, from \eqref{eq-omega(t)-bound}, 
\eqref{eq-slowdown} is immediate. 
% \begin{equation}
%     \lambda_{1,2} = 1+\frac{c_1^2}{2} \pm \frac{c_1\sqrt{c_1^2 + 4}}{2}
% \end{equation}
\end{proof}

\end{document}
